% Source: https://artofproblemsolving.com/community/c3h2302883_solution_without_sines
% Stable identity: aops-topic-2302883-q01
The pole, the vertical drop and the ground form a right triangle with hypotenuse $5$ m and an angle of $20^\circ$ at the base of the pole. The honey lands at the foot of the vertical through the top, at distance $d$ from the base:

\begin{asy}
usepackage("amsmath");
size(6cm);
defaultpen(fontsize(10pt));
pair O=(0,0), T=5*dir(20), F=(T.x,0);
draw((-0.6,0)--(5.6,0),linewidth(0.8));
draw(O--T,brown+linewidth(2));
draw(T--F,dashed+orange,Arrow(5));
draw(F+(-0.15,0)--F+(-0.15,0.15)--F+(0,0.15));
filldraw(circle(T,0.09),orange,orange);
draw(arc(O,0.8,0,20)); label("$20^\circ$",0.85*dir(8),E,fontsize(8pt));
label("$5$ m",(O+T)/2,NW); label("$d$",(O+F)/2,S,blue);
label("base",O,SW,fontsize(8pt)); label("honey",T,N,fontsize(8pt));
\end{asy}

\[d=5\cos20^\circ\approx5\cdot0.93969\approx\boxed{4.70\text{ m}}.\]
A trigonometric value cannot be avoided here, because $20^\circ$ is not one of the angles of the special right triangles ($30^\circ$, $45^\circ$, $60^\circ$); one needs a table, a calculator, an accurate scale drawing, or some equation that $\cos20^\circ$ satisfies. (For example, the triple-angle identity $\cos60^\circ=4\cos^320^\circ-3\cos20^\circ$ shows that $d$ is the root near $4.7$ of $8d^3-150d-125=0$, which can be found to any accuracy by trial.)
